% SPDX-License-Identifier: MIT
% Copyright (c) 2026 Evgeny Baulin
% Optimization Methods, Seminar 10: one-page cheat sheet (Russian).
% Generated by checks/10. Knapsack and Linear Programming in Practice.ipynb from web/10/js/content-cheatsheet.js and content/students.json. Do not edit by hand.
% Compile with: latexmk -xelatex -interaction=nonstopmode -halt-on-error Cheatsheet_ru.tex
\RequirePackage{fix-cm}
\documentclass[10pt,a4paper]{article}

\usepackage{amsmath,amssymb}
\usepackage{fontspec}
\setmainfont{cmunrm.otf}[BoldFont=cmunbx.otf, ItalicFont=cmunti.otf, BoldItalicFont=cmunbi.otf]
\setsansfont{cmunss.otf}[BoldFont=cmunsx.otf, ItalicFont=cmunsi.otf, BoldItalicFont=cmunso.otf]
\setmonofont{cmuntt.otf}[BoldFont=cmuntb.otf, ItalicFont=cmunit.otf, BoldItalicFont=cmuntx.otf]
\usepackage{polyglossia}
\setdefaultlanguage{russian}
\setotherlanguage{english}
\newfontfamily\cyrillicfont{cmunrm.otf}[Ligatures=TeX, BoldFont=cmunbx.otf, ItalicFont=cmunti.otf, BoldItalicFont=cmunbi.otf]
\newfontfamily\cyrillicfontsf{cmunss.otf}[Ligatures=TeX, BoldFont=cmunsx.otf, ItalicFont=cmunsi.otf, BoldItalicFont=cmunso.otf]
\newfontfamily\cyrillicfonttt{cmuntt.otf}[BoldFont=cmuntb.otf, ItalicFont=cmunit.otf, BoldItalicFont=cmuntx.otf]
\PolyglossiaSetup{russian}{hyphenmins={2,3}}

\usepackage[a4paper,left=8mm,right=8mm,top=7mm,bottom=11mm,footskip=5.5mm]{geometry}
\usepackage{xcolor}
\usepackage{multicol}
\usepackage{enumitem}
\usepackage{ragged2e}
\usepackage{fancyhdr}
\usepackage[colorlinks=true,linkcolor=hseblue,urlcolor=hseblue]{hyperref}

\definecolor{hseblue}{RGB}{16,45,105}
\definecolor{hseaccent}{RGB}{198,40,40}
\definecolor{hsegray}{RGB}{108,114,126}
\definecolor{hselight}{RGB}{240,243,248}

\hypersetup{pdftitle={Шпаргалка к семинару 10: Рюкзак, ветви и границы, таблицы и модели}, pdfauthor={Евгений Баулин}}

\pagestyle{fancy}
\fancyhf{}
\renewcommand{\headrulewidth}{0pt}
\fancyfoot[L]{\scriptsize\color{hsegray}\textcopyright\ 2026 Евгений Баулин, \href{https://t.me/tarakan_tuc}{t.me/tarakan\_tuc}}
\fancyfoot[R]{\scriptsize\color{hsegray}Лицензия MIT}

\setlength{\parindent}{0pt}
\setlength{\parskip}{0pt}
\setlength{\columnsep}{4.5mm}
\setlength{\columnseprule}{0.3pt}
\renewcommand{\columnseprulecolor}{\color{hselight}}
\setlength{\emergencystretch}{2em}
\setlength{\RaggedRightRightskip}{0pt plus 4em}
\binoppenalty=2000
\relpenalty=500
\setlength{\abovedisplayskip}{2pt}
\setlength{\belowdisplayskip}{2pt}
\setlist[itemize]{leftmargin=7pt, labelsep=3pt, itemsep=1.2pt, topsep=1pt, parsep=0pt, partopsep=0pt,
  label={\color{hseaccent}\scriptsize\textbullet}}

\makeatletter
\newcommand{\cssec}[1]{\par\addvspace{3.5pt}\noindent
  \colorbox{hselight}{\parbox{\dimexpr\linewidth-2\fboxsep}{\RaggedRight\hyphenpenalty=10000\color{hseblue}\bfseries\strut #1}}%
  \par\nobreak\vspace{1pt}\@afterheading}
\makeatother
\clubpenalty=10000
\widowpenalty=10000

\begin{document}
\thispagestyle{fancy}

\noindent
\begin{minipage}[b]{0.6\textwidth}
  \RaggedRight\hyphenpenalty=10000
  {\color{hseaccent}\bfseries\small Шпаргалка к семинару 10}\\[1pt]
  {\color{hseblue}\bfseries\Large Рюкзак, ветви и границы,\\ таблицы и модели}
\end{minipage}\hfill
\begin{minipage}[b]{0.39\textwidth}
  \raggedleft\footnotesize
  {\bfseries Методы оптимизации}\\
  НИУ ВШЭ, факультет компьютерных наук, 2026--2027\\
  Лектор курса: Андрей Игнатов\\
  Автор: Евгений Баулин, Telegram \href{https://t.me/tarakan_tuc}{@tarakan\_tuc}
\end{minipage}

\vspace{2pt}
{\color{hseblue}\rule{\textwidth}{0.8pt}}
\vspace{-4pt}

\fontsize{9}{10.6}\selectfont
\begin{multicols}{3}
\RaggedRight


% ---- knap
\cssec{Рюкзак}
\begin{itemize}
  \item $\max p^{\top}x$, $w^{\top}x\le W$, $x\in\{0,1\}^n$; релаксация $0\le x\le1$.
  \item Релаксация: сортировка по $p_i/w_i$, целые предметы, одна доля; значение $U$.
  \item Целые ценности: $\text{opt}\le\lfloor U\rfloor$. Округление вниз — лишь эвристика.
\end{itemize}

% ---- bnb
\cssec{Ветви и границы}
\begin{itemize}
  \item Узел $N$: зафиксированные переменные; $U(N)$ = зафиксированная ценность + релаксация остального.
  \item Закрыть, если недопустим, если $U(N)\le z_{\text{inc}}$ или если целый (обновить $z_{\text{inc}}$).
  \item Ветвиться по дробной переменной: $x_j=1$, $x_j=0$.
  \item В глубину: память, ранний рекорд. По лучшей оценке: меньше узлов.
  \item В любой момент: $z_{\text{inc}}\le\text{opt}\le\max\bigl(z_{\text{inc}},\max_{\text{open}}U(N)\bigr)$.
\end{itemize}

% ---- cut
\cssec{Отсечение Гомори}
\begin{itemize}
  \item Строка $x_k+\sum_N\bar a_jx_j=\bar b$, $\bar b\notin\mathbb Z$: $\sum_N\{\bar a_j\}x_j\ge\{\bar b\}$.
  \item Верно для целых точек, нарушено оптимумом ЛП; подставьте дополнительные переменные, чтобы записать в $x$.
\end{itemize}

% ---- dp
\cssec{Таблица ДП}
\begin{itemize}
  \item $D(i,c)=\max\{D(i-1,c),\,p_i+D(i-1,c-w_i)\}$, $D(0,c)=0$, ответ $D(n,W)$.
  \item Назад: $D(i,c)=D(i-1,c)\Rightarrow$ не взят; иначе взят, $c\mathrel{-}=w_i$.
  \item Одна строка: для $0$–$1$ проход по $c$ вниз; неограниченные копии — вверх.
  \item $O(nW)$ — псевдополиномиально; сумма подмножества — случай $p=w$.
\end{itemize}

% ---- model
\cssec{Двоичные переменные}
\begin{itemize}
  \item $A\Rightarrow B$: $y_A\le y_B$; не больше $k$: $\sum y\le k$.
  \item Фиксированные затраты: $x\le u\,y$, $-f\,y$; наименьшая верная $u$.
  \item $z=x_ix_j$: $z\le x_i$, $z\le x_j$, $z\ge x_i+x_j-1$.
  \item $z=y\,x$, $0\le x\le u$: $z\le uy$, $z\le x$, $z\ge x-u(1-y)$.
\end{itemize}

% ---- forms
\cssec{Модели}
\begin{itemize}
  \item $P_1\subseteq P_2$: $P_1$ сильнее, оценка лучше, дерево меньше.
  \item Разделяйте: $x_{ij}\le y_j$ лучше $\sum_ix_{ij}\le m\,y_j$; клики лучше пар.
\end{itemize}

\end{multicols}
\end{document}
